% T19 -- AOP (amplificateur op\'erationnel) -- banque d'exercices
% Ordre des exercices : gradation croissante 1,1,2,2,2,3
% Principe : banque COMPLEMENTAIRE de la fiche T19_AOP.tex -- aucun exercice ne
% reprend les donnees ni le travail des Questions Types QT1 (2024 E1-E3, G1=3,4,
% gain global 340), QT2 (2023 F6-F8, seuils 2,50/2,30 V) et QT3 (2024 E9, Is=V4/R5).
% Sources :
%   - Exercice 1 : E2_Sujet_2026_BAC3 (questions B4 et B5) : conditionneur de sonde de
%     temp\'erature Jumo par source de courant constant I=10 mA et montage suiveur, RT=R0(1+aT).
%     Cahier des charges du sujet (avant Q.B1) : \'etendue de mesure -5 \`a 90 degC,
%     erreur de mesure < 1 degC. Q.B5 demande vT aux DEUX temp\'eratures extr\^emes du CDC
%     (-5 et 90 degC). Q.B6 : la plage de tension exig\'ee en entr\'ee du CAN est de 3 V.
%     R0=100 Ohm et a=3,85e-3 /K sont donn\'es ici comme valeurs typiques (le sujet les fait
%     extraire du document technique DTB2, non reproduit) : aucune attribution au sujet.
%   - Exercice 2 (docTech) : E1_Sujet_2024 (= E2_Disciplinaire_IE_2024, doublon de conversion ;
%     \'epreuve \'ecrite disciplinaire ; extracteur d'air), question E5 + document technique DT10.
%     TOUT le contenu de l'encadr\'e docTech est TEXTUELLEMENT pr\'esent dans le .md du sujet :
%     "thermocouple type J", plage -20 \`a 200 degC, tableau des points caract\'eristiques
%     (eABmini = -1 mV, eABmaxi = 10,8 mV), profil du courant Is mini = 4 mA / Is maxi = 20 mA,
%     raccordement 2 fils (capteur) et 4 fils (sortie), lin\'earit\'e des profils (\'enonc\'ee
%     avant Q.E1). La relation V4 = 337,7*eAB + 1,34 est donn\'ee par le sujet lui-m\^eme
%     (juste avant E5). RECONSTITUE : seule la mise en forme (tableau LaTeX au lieu de la
%     figure) -- l'image du DT10 n'est pas restitu\'ee par l'OCR. Les r\'esultats V4 = 1 V et
%     V4 = 5 V sont confirm\'es par le Rapport_Jury_2024 (corrig\'e officiel E5).
%   - Exercice 3 : m\^eme sujet 2024, question E4 seule : \'etage de d\'ecalage en tension,
%     th\'eor\`eme de Millman, V4 = V3 + Ve -- relations de Millman donn\'ees par la figure 23
%     du sujet (texte OCR restitu\'e).
%   - Exercice 4 (tableau comparatif des montages) : exercice in\'edit de synth\`ese.
%   - Exercice 5 : E2_Sujet_2026_BAC3, questions B7 et B8 + document
%     technique DTB3. Les quatre montages et leurs fonctions de transfert sont TEXTUELLEMENT
%     pr\'esents dans le .md (bloc "Document technique DTB3", montages 1 \`a 4 : vS/vE=(R1+R0)/R1,
%     vS/vE=-R0/R1, vS=(R0/R1)(v2-v1), et montage 4 \`a R1/R2/R3/R0). RECONSTITUE : la colonne
%     "c\^ablage", d\'eduite sans ambigu\"it\'e de chaque fonction de transfert (les figures ne
%     sont pas restitu\'ees par l'OCR). Le sujet indique lui-m\^eme (avant B9) que "le montage 3
%     est finalement retenu", ce qui valide la r\'eponse attendue en B8.
%   - Exercice 6 : E2_Disciplinaire_IE_2023 (robot Bakus, questions F9 et F10), pris ici sous
%     l'angle du DIMENSIONNEMENT (probl\`eme inverse) et non du calcul direct des seuils, qui
%     est le travail de la QT2 de la fiche. Donn\'ees du sujet : sorties inertielles 0-5 V pour
%     0-50 % de pente ; basculement autonome->manuel au-dessus de 25 % ; retour au mode
%     autonome seulement en dessous de 23 % ; ALI aliment\'e en 15 V, tensions de saturation
%     15 V ; valeurs retenues par le concepteur R1=1 kOhm, R2=150 kOhm, Vref2=2,42 V.
%     Figure 23 du sujet : Vref2 est reli\'ee au noeud V+ par R1 (1 kOhm) et la sortie Vs2 par
%     R2 (150 kOhm) -- orientation conserv\'ee du lot 2, coh\'erente avec le corrig\'e officiel
%     F7/F8 (V+ = 378/151 = 2,5 V et 348/151 = 2,3 V) et avec F10 (conformit\'e, retour au
%     mode autonome en dessous de 23 % de pente). R2 = 149 kOhm et Vref2 = 2,416 V sont les
%     valeurs THEORIQUES recalcul\'ees ici \`a partir des seuils 2,5/2,3 V du cahier des
%     charges ; elles ne sont pas attribu\'ees au sujet, qui ne donne que les valeurs
%     normalis\'ees finales.
\expandafter\ifx\csname CAPETINCLUDEMODE\endcsname\relax
\documentclass[11pt,a4paper]{article}
% preamble.tex — packages et commandes communes a toutes les fiches CAPET SII IE
\input{glyphtounicode}
\pdfgentounicode=1
\usepackage{cmap}
\usepackage[T1]{fontenc}
\usepackage[utf8]{inputenc}
\usepackage[french]{babel}
% babel[french] active a \begin{document} des raccourcis de ponctuation (:;!?)
% qui cassent la syntaxe to[...] de circuitikz (erreur "+ or - expected") ;
% on les desactive, ce qui ne change que l'espacement fin avant ces signes.
\AtBeginDocument{\shorthandoff{:;!?}}
\usepackage[a4paper, top=1.6cm, bottom=1.6cm, left=1.4cm, right=1.4cm]{geometry}
\setlength{\headheight}{14.5pt}
\usepackage{amsmath, amssymb, amsfonts}
\usepackage{siunitx}
\sisetup{locale=FR, output-decimal-marker={,}, inter-unit-product=\ensuremath{{}\cdot{}},
         group-minimum-digits=4}
% Unites non fournies par siunitx, declarees ici une fois pour toutes (fiches + exercices).
% Les gardes \ifdefined evitent tout conflit avec une redeclaration locale restee dans un
% fichier : ne pas retirer les gardes, ne pas redeclarer ces unites dans les fichiers.
\DeclareSIUnit{\jour}{j}
\ifdefined\var\else\DeclareSIUnit{\var}{var}\fi   % puissance reactive
\ifdefined\vad\else\DeclareSIUnit{\vad}{VAD}\fi   % puissance deformante
\ifdefined\tr\else\DeclareSIUnit{\tr}{tr}\fi      % tour (tr/min)
\usepackage{array, booktabs, multirow, tabularx}
\usepackage{xcolor}
\usepackage[most, breakable]{tcolorbox}
\usepackage{multicol}
\usepackage{enumitem}
\usepackage{mathtools}
\usepackage{lmodern}
% Refonte visuelle 2026 (direction "moderne coloré") : Lato pour le texte,
% maths sans-serif assorties via newtxsf (charge apres amsmath/amssymb).
\usepackage[default]{lato}
\usepackage{newtxsf}
\usepackage{microtype}
\usepackage{fancyhdr}
\usepackage{tikz}
\usetikzlibrary{positioning,arrows.meta}
\usepackage{circuitikz}
\usepackage{lastpage}

% ── Couleurs par domaine (palette modernisée 2026) ────────────────────
\definecolor{colElec}{RGB}{29, 78, 216}        % Électrotechnique — bleu vif
\definecolor{colElecLight}{RGB}{219, 234, 254} % fond bleu clair
\definecolor{colAuto}{RGB}{21, 128, 61}        % Automatique — vert
\definecolor{colMeca}{RGB}{194, 65, 12}        % Mécanique — orange
\definecolor{colTherm}{RGB}{185, 28, 28}       % Thermique — rouge
\definecolor{colQT}{RGB}{232, 244, 233}        % (hérité, plus utilisé par questiontype)
\definecolor{colJury}{RGB}{255, 248, 225}      % fond jury (crème)
\definecolor{colJuryFrame}{RGB}{202, 138, 4}   % bordure/bandeau jury (ambre)
\definecolor{colGrayLight}{RGB}{245, 245, 247} % fond gris neutre
\definecolor{colSeparator}{RGB}{180, 180, 180} % gris séparateur
\definecolor{colExo}{RGB}{8, 145, 178}          % Banque d'exercices — cyan (distinct des fiches)
% ── Encre et gris de la refonte ───────────────────────────────────────
\definecolor{colInk}{RGB}{23, 32, 42}          % texte courant (quasi-noir)
\definecolor{colSlate}{RGB}{83, 98, 118}       % gris ardoise (méta, en-têtes)
\definecolor{colRule}{RGB}{203, 213, 225}      % filets clairs
\definecolor{colJuryBg}{RGB}{255, 248, 225}    % fond crème jury
\definecolor{colPresence}{RGB}{83, 98, 118}    % bloc 5 (présence dans les sujets)
\AtBeginDocument{\color{colInk}}

% ── Style encadré exercice + corrigé (banque d'exercices) ─────────────
\tcbset{
  exoBox/.style={
    breakable,
    enhanced,
    colback=colExo!8!white,
    colframe=colExo,
    colbacktitle=colExo,
    coltitle=white,
    fonttitle=\bfseries\small,
    left=4pt, right=4pt, top=3pt, bottom=3pt,
    before skip=5pt, after skip=5pt,
  }
}

% ── Bascule enonce / corrige (banque d'exercices) ─────────────────────
% \CAPETCORRIGE est defini en ligne de commande (pdflatex "\def\CAPETCORRIGE{1}\input{...}")
% pour compiler la variante corriges d'un meme fichier source ; par defaut (non defini),
% la compilation produit la variante enonces. \CAPETSHOW{enonce}{corrige} choisit la branche.
%
% TROIS MODES (refonte 2026-09-20, manuel de revision) :
%   - aucun drapeau   -> imprime l'enonce            (variante enonces, INCHANGE)
%   - \CAPETCORRIGE   -> imprime le corrige          (variante corriges, INCHANGE)
%   - \CAPETMANUEL    -> imprime l'enonce ET met le corrige de cote pour l'annexe
%
% Le mode manuel stocke chaque corrige dans une macro (\csgdef), sans ecriture
% fichier ni re-tokenisation : le contenu (circuitikz, accents, %) traverse intact.
% L'annexe les rejoue ensuite dans l'ordre via \CAPETANNEXE.
\usepackage{etoolbox}

\newcounter{CAPETexo}          % numero continu d'exercice dans le manuel
\setcounter{CAPETexo}{0}
\newcounter{CAPETannexe}       % compteur de relecture en annexe

% Titre courant de la notion en cours (pose par \CAPETNOTION), sert d'intertitre annexe
\newcommand{\CAPETnotioncode}{}
\newcommand{\CAPETnotiontitre}{}

% \CAPETNOTION{code}{titre} : ouvre une notion dans le manuel (memorise le libelle,
% et marque le rang du premier exercice de la notion pour l'intertitre d'annexe).
\newcommand{\CAPETNOTION}[2]{%
  \renewcommand{\CAPETnotioncode}{#1}%
  \renewcommand{\CAPETnotiontitre}{#2}%
  \csgdef{CAPETnotionstart\the\numexpr\value{CAPETexo}+1\relax}{#1 -- #2}%
}

% ── \degree robuste ───────────────────────────────────────────────────
% Certains fichiers d'exercices utilisent \degree. Ni le preambule ni siunitx
% ne le definissent, et \providecommand local ne suffit pas : sous hyperref,
% \degree existe deja comme commande factice (\ERROR), donc \providecommand
% s'abstient et l'appel echoue a la compilation du manuel (erreur invisible
% en compilation isolee). On le (re)definit donc inconditionnellement ici,
% apres le chargement des packages, pour tout le corpus.
\AtBeginDocument{\providecommand{\degree}{}\renewcommand{\degree}{\ensuremath{^\circ}}}

% Libelle decoratif : les fichiers d'exercices utilisent aussi \CAPETSHOW dans
% leur \fancyhead et leur bloc de titre, pour afficher "(enonces)" ou "(corriges)".
% Ces appels-la ne sont PAS des exercices : en mode manuel ils ne doivent ni
% consommer un numero, ni alimenter l'annexe. On les reconnait a leur premier
% argument, strictement egal au libelle ci-dessous.
\newcommand{\CAPETLIBELLEENONCES}{(\'enonc\'es)}

\ifdefined\CAPETMANUEL
  % --- Mode manuel : enonce en place, corrige differe vers l'annexe ---
  \newcommand{\CAPETSHOW}[2]{%
    \def\CAPETargun{#1}%
    \ifx\CAPETargun\CAPETLIBELLEENONCES
      #1%            appel decoratif : rendu tel quel, hors numerotation
    \else
      \CAPETSHOWEXO{#1}{#2}%
    \fi
  }
  \newcommand{\CAPETSHOWEXO}[2]{%
    \stepcounter{CAPETexo}%
    \csgdef{CAPETcorrige\theCAPETexo}{#2}%
    \edef\CAPETtmpnum{\theCAPETexo}%
    #1%
    \par\vspace{2pt}%
    {\footnotesize\itshape\color{colExo}%
     \hyperref[capet:cor:\CAPETtmpnum]{Corrige de l'exercice \CAPETtmpnum\ $\rightarrow$ p.~\pageref{capet:cor:\CAPETtmpnum}}%
     \label{capet:eno:\CAPETtmpnum}}%
  }
  % \CAPETANNEXE : rejoue tous les corriges captures, dans l'ordre, avec intertitres
  \newcommand{\CAPETANNEXE}{%
    \setcounter{CAPETannexe}{0}%
    \loop\ifnum\value{CAPETannexe}<\value{CAPETexo}%
      \stepcounter{CAPETannexe}%
      \edef\CAPETcurnum{\theCAPETannexe}%
      \ifcsname CAPETnotionstart\CAPETcurnum\endcsname
        \par\medskip
        {\normalsize\bfseries\color{colExo}\csuse{CAPETnotionstart\CAPETcurnum}}%
        \par\nobreak\smallskip
      \fi
      \begin{tcolorbox}[exoBox, breakable,
        title={Corrige de l'exercice \CAPETcurnum\
               \hfill\normalfont\footnotesize
               \hyperref[capet:eno:\CAPETcurnum]{enonce p.~\pageref{capet:eno:\CAPETcurnum}}}]%
      \label{capet:cor:\CAPETcurnum}%
      \csuse{CAPETcorrige\CAPETcurnum}%
      \end{tcolorbox}%
    \repeat
  }
\else
  % --- Modes historiques, comportement strictement inchange ---
  \newcommand{\CAPETSHOW}[2]{\ifdefined\CAPETCORRIGE #2\else #1\fi}
  \newcommand{\CAPETANNEXE}{}
\fi

% ── Difficulté 1 à 3 (étoiles pleines, amssymb) ───────────────────────
\newcommand{\difficulte}[1]{%
  \ifcase#1\relax
  \or $\bigstar$\;niveau 1/3%
  \or $\bigstar\bigstar$\;niveau 2/3%
  \or $\bigstar\bigstar\bigstar$\;niveau 3/3%
  \fi
}

% ── Style encadré extrait de document technique (banque d'exercices) ──
\tcbset{
  docTech/.style={
    breakable,
    enhanced,
    colback=colGrayLight,
    colframe=colSeparator,
    fonttitle=\bfseries\small\itshape,
    coltitle=black,
    colbacktitle=colGrayLight,
    left=4pt, right=4pt, top=3pt, bottom=3pt,
    before skip=4pt, after skip=4pt,
  }
}

% ── Style tcolorbox de section principale ────────────────────────────
% Refonte 2026 : bandeau titre coloré arrondi + fond très pâle, sans cadre.
\tcbset{
  ficheSection/.style={
    breakable, enhanced jigsaw,
    frame hidden, boxrule=0pt, arc=1.8mm,
    colback=#1!4!white,
    colbacktitle=#1,
    coltitle=white,
    fonttitle=\sffamily\bfseries,
    left=7pt, right=7pt, top=5pt, bottom=5pt,
    toptitle=3pt, bottomtitle=3pt,
    before skip=10pt, after skip=8pt,
  }
}

% ── Style encadré question type ──────────────────────────────────────
% Carte blanche, bandeau sombre, liseré très clair.
\tcbset{
  questiontype/.style={
    breakable, enhanced jigsaw,
    colback=white,
    frame hidden, boxrule=0pt, arc=1.8mm,
    colbacktitle=colElec!78!black,
    coltitle=white,
    fonttitle=\sffamily\bfseries\small,
    left=7pt, right=7pt, top=4pt, bottom=4pt,
    toptitle=2.5pt, bottomtitle=2.5pt,
    before skip=8pt, after skip=8pt,
    borderline={0.6pt}{0pt}{colElec!12!white},
  }
}

% ── Style encadré jury ───────────────────────────────────────────────
% Bandeau ambre, fond crème, sans cadre.
\tcbset{
  jury/.style={
    breakable, enhanced jigsaw,
    colback=colJuryBg,
    frame hidden, boxrule=0pt, arc=1.8mm,
    colbacktitle=colJuryFrame,
    coltitle=white,
    fonttitle=\sffamily\bfseries,
    left=7pt, right=7pt, top=4pt, bottom=4pt,
    toptitle=3pt, bottomtitle=3pt,
    before skip=10pt, after skip=8pt,
  }
}

% ── En-tête / pied de page communs (refonte 2026) ─────────────────────
% #1 code fiche, #2 titre court, #3 couleur domaine
\newcommand{\ficheHeaders}[3]{%
  \fancyhead[L]{\small\sffamily{\bfseries\textcolor{#3}{#1}}\hspace{0.7em}\textcolor{colSlate}{#2}}%
  \fancyhead[R]{\small\sffamily\textcolor{colSlate}{CAPET SII IE\, --- \,Session 2026}}%
  \fancyfoot[C]{\small\sffamily\textcolor{colSlate}{\thepage\,/\,\pageref{LastPage}}}%
  \fancyhead[C]{}\fancyfoot[L]{}\fancyfoot[R]{}%
}

% ── Bloc en-tête de fiche (BLOC 0, refonte 2026) ──────────────────────
% #1 code  #2 titre  #3 couleur domaine  #4 domaine  #5 ligne présence  #6 sous-thèmes
\newcommand{\ficheEntete}[6]{%
  \begin{tcolorbox}[enhanced, breakable,
    frame hidden, boxrule=0pt, arc=2.5mm,
    interior style={left color=#3, right color=#3!75!black},
    left=9pt, right=9pt, top=7pt, bottom=7pt,
    before skip=0pt, after skip=6pt]
    {\color{white}\sffamily\bfseries\LARGE #1\quad #2}\\[4pt]
    {\color{#3!22!white}\sffamily\small \MakeUppercase{#4}\quad\textbar\quad #5}\\[1pt]
    {\color{#3!35!white}\sffamily\footnotesize Sous-thèmes : #6}
  \end{tcolorbox}%
}

% ── Mise en page ─────────────────────────────────────────────────────
\pagestyle{fancy}
\fancyhf{}
\renewcommand{\headrulewidth}{0.5pt}
\renewcommand{\footrulewidth}{0pt}
\AtBeginDocument{%
  \renewcommand{\headrule}{{\color{colRule}\hrule width\headwidth height 0.5pt}}%
}

\setlist[itemize]{leftmargin=1.2em, itemsep=1pt, topsep=2pt, parsep=0pt}
\setlist[enumerate]{leftmargin=1.2em, itemsep=3pt, topsep=2pt, parsep=0pt}

\setlength{\parskip}{2pt}
\setlength{\columnsep}{10pt}

\def\CAPETBEGINFILE{\begin{document}}
\def\CAPETENDFILE{\end{document}}
\else
\def\CAPETBEGINFILE{}
\def\CAPETENDFILE{}
\fi

\fancyhead[L]{\textbf{\textcolor{colExo}{T19}} - Exercices - AOP (amplificateur op\'erationnel)}
\fancyhead[R]{\small\textit{Banque d'exercices \CAPETSHOW{(\'enonc\'es)}{(corrig\'es)}}}
\fancyfoot[C]{\small\thepage\,/\,\pageref{LastPage}}

\CAPETBEGINFILE

\begin{tcolorbox}[
  enhanced, breakable,
  colback=colExo,
  colframe=colExo,
  left=6pt, right=6pt, top=5pt, bottom=5pt,
  before skip=0pt, after skip=4pt
]
  {\color{white}\bfseries\large T19 \quad AOP (amplificateur op\'erationnel)}\\[2pt]
  {\color{white}\small Banque d'exercices -- 6 exercices \CAPETSHOW{(\'enonc\'es)}{(corrig\'es)}}
\end{tcolorbox}

\begin{tcolorbox}[exoBox,
  title={Exercice 1 -- \difficulte{1} \hfill\normalfont\footnotesize 2026, \'epreuve BAC3, questions B4 et B5}
]
\CAPETSHOW{
Un capteur de temp\'erature Jumo est mod\'elis\'e par une r\'esistance $R_T = R_0(1+aT)$, aliment\'ee par une source de courant constant $I = \SI{10}{\milli\ampere}$. La tension aux bornes de $R_T$ attaque l'entr\'ee non-inverseuse d'un AOP suppos\'e id\'eal, c\^abl\'e comme indiqu\'e ci-dessous~; sa sortie $v_T$ est transmise aux \'etages suivants. Le cahier des charges de la mesure impose une \'etendue de mesure de $\SI{-5}{\celsius}$ \`a $\SI{90}{\celsius}$ et une erreur de mesure inf\'erieure \`a $\SI{1}{\celsius}$.

\begin{center}
\begin{circuitikz}[european, scale=0.85]
\coordinate (LO) at (0,2.05);
\draw (0,0) to[I, l=$I$] (0,2.5) to[R, l=$R_T$] (0,5) coordinate(T);
\draw (0,0) node[ground]{};
\draw (4,3.75) node[op amp, yscale=-1] (U) {};
\draw (T) -- (2.2,5) |- (U.+);
\draw (U.out) -- ++(1.2,0) coordinate(S) node[circ]{} -- ++(0.9,0) node[right]{$v_T$};
\draw (U.-) -- ++(-1.5,0) coordinate(Nn) -- (Nn |- LO) -- (S |- LO) -- (S);
\end{circuitikz}
\end{center}

\begin{enumerate}[leftmargin=1.5em]
  \item Exprimer $v_T$ en fonction de $I$, $R_0$ et $a\cdot T$, sous la forme $v_T = U_0(1+aT)$ en pr\'ecisant $U_0$.
  \item Quel est le nom de ce montage \`a AOP, et quel est son r\^ole fonctionnel dans cette cha\^ine de conditionnement~?
  \item Sachant que $R_0 = \SI{100}{\ohm}$ et $a = \SI{3,85e-3}{\per\kelvin}$ (valeurs typiques d'une sonde platine), calculer $v_T$ pour les deux temp\'eratures extr\^emes du cahier des charges.
\end{enumerate}

}{
\textbf{Corrig\'e.}

\begin{enumerate}[leftmargin=1.5em]
  \item La source de courant $I$ impose $I$ dans $R_T$ (l'AOP suiveur ne pr\'el\`eve aucun courant sur son entr\'ee $+$, donc tout $I$ traverse bien $R_T$) : par la loi d'Ohm,
  \[ v_T = R_T\cdot I = I\cdot R_0(1+aT) = \underbrace{I\cdot R_0}_{U_0}\,(1+aT) \]
  soit $U_0 = I\cdot R_0$.
  \item C'est un montage \textbf{suiveur} (buffer, gain unitaire $v_T=V_+$). Son r\^ole est l'\textbf{adaptation d'imp\'edance} : il isole le capteur $R_T$ (imp\'edance de sortie non nulle) de l'\'etage suivant, en pr\'esentant une imp\'edance d'entr\'ee quasi infinie \`a $R_T$ (aucun courant pr\'elev\'e, la mesure n'est donc pas perturb\'ee) et une imp\'edance de sortie quasi nulle pour attaquer la suite de la cha\^ine sans chute de tension suppl\'ementaire.
  \item $U_0 = I\cdot R_0 = 0{,}010\times 100 = \SI{1}{\volt}$.

  Pour $T=\SI{-5}{\celsius}$ : $v_T = 1\times(1-3{,}85\times10^{-3}\times 5) = 1\times(1-0{,}01925) \approx \SI{0,981}{\volt}$.

  Pour $T=\SI{90}{\celsius}$ : $v_T = 1\times(1+3{,}85\times10^{-3}\times 90) = 1\times(1+0{,}3465) \approx \SI{1,347}{\volt}$.
\end{enumerate}

\textbf{Conclusion.} Sur l'\'etendue de mesure $-5$ \`a $\SI{90}{\celsius}$ fix\'ee par le cahier des charges, la tension image de la temp\'erature ne balaie que $\SI{0,981}{\volt}$ \`a $\SI{1,347}{\volt}$, soit une sensibilit\'e $U_0\,a = \SI{3,85}{\milli\volt\per\kelvin}$ et une dynamique utile de $\SI{0,366}{\volt}$ seulement. Le suiveur garantit que cette faible tension est pr\'elev\'ee \textbf{sans charger la sonde} -- donc sans erreur syst\'ematique de mesure, condition n\'ecessaire pour tenir l'exigence d'erreur inf\'erieure \`a $\SI{1}{\celsius}$ --, mais il n'apporte aucun gain : la plage de $\SI{3}{\volt}$ exig\'ee en entr\'ee du convertisseur analogique-num\'erique impose donc, en aval, un \'etage amplificateur-conditionneur.
}
\end{tcolorbox}

\begin{tcolorbox}[exoBox,
  title={Exercice 2 -- \difficulte{1} \hfill\normalfont\footnotesize 2024, \'epreuve \'ecrite disciplinaire, question E5 (document technique DT10)}
]
\CAPETSHOW{
Un transmetteur de temp\'erature convertit la tension $e_{AB}$ d\'elivr\'ee par un capteur en un courant normalis\'e $I_s$ exploitable par l'entr\'ee analogique d'un automate. Il encha\^ine trois fonctions : \emph{amplifier} ($\times G$), \emph{d\'ecaler} ($+V_d$) puis \emph{convertir} la tension en courant. \`A l'issue des deux premi\`eres, la tension $V_4$ v\'erifie la relation, donn\'ee par le sujet :
\[ V_4 = 337{,}7\ e_{AB} + 1{,}34 \qquad (e_{AB}\ \text{en volts}) \]

\begin{tcolorbox}[docTech, title={Document technique DT10 -- Profil de la tension fournie par le capteur de temp\'erature}]
\footnotesize
Capteur de temp\'erature de type \textbf{thermocouple type J}, \'etalonn\'e sur la plage $-20$ \`a $\SI{200}{\celsius}$, raccord\'e 2 fils. Le profil $e_{AB}(T_{ARM})$ est \textbf{lin\'eaire} entre les deux points caract\'eristiques ci-dessous.

\smallskip
\begin{center}
\renewcommand{\arraystretch}{1.2}
\begin{tabular}{lcc}
\toprule
\textbf{Points caract\'eristiques maximaux} & & \\
\midrule
Temp\'erature mesur\'ee ($\si{\celsius}$) & $-20$ & $200$ \\
Tension d\'elivr\'ee par le capteur ($\si{\milli\volt}$) & $e_{AB\,mini} = -1$ & $e_{AB\,maxi} = 10{,}8$ \\
\bottomrule
\end{tabular}
\end{center}

\smallskip
Profil du courant d\'elivr\'e par le transmetteur $TT$ : lin\'eaire \'egalement, de $I_{s\,mini} = \SI{4}{\milli\ampere}$ (\`a $\SI{-20}{\celsius}$) \`a $I_{s\,maxi} = \SI{20}{\milli\ampere}$ (\`a $\SI{200}{\celsius}$), sortie 4 fils.
\end{tcolorbox}

\begin{enumerate}[leftmargin=1.5em]
  \item \textbf{Relever} dans le document technique DT10 : la nature du transducteur, l'\'etendue de mesure en temp\'erature, et les deux valeurs extr\^emes de la tension $e_{AB}$ d\'elivr\'ee par le capteur.
  \item Calculer la tension $V_4$ pour chacune de ces deux tensions extr\^emes de $e_{AB}$.
  \item En d\'eduire la sensibilit\'e $\Delta V_4/\Delta T_{ARM}$ de la cha\^ine \{capteur + amplification + d\'ecalage\}, exprim\'ee en $\si{\milli\volt\per\celsius}$.
  \item Quel est le r\^ole du terme constant $1{,}34$ dans la relation $V_4 = 337{,}7\,e_{AB} + 1{,}34$~? Que vaudrait $V_4$ \`a $\SI{-20}{\celsius}$ sans ce terme, et pourquoi serait-ce r\'edhibitoire pour la suite de la cha\^ine~?
\end{enumerate}

}{
\textbf{Corrig\'e.}

\begin{enumerate}[leftmargin=1.5em]
  \item \textbf{Extraction du DT10.}
  \begin{itemize}
    \item Transducteur : \textbf{thermocouple de type J} -- capteur \textbf{actif} (il g\'en\`ere lui-m\^eme une f.\'e.m. par effet Seebeck, sans alimentation externe), ce qui explique que la tension utile ne soit que de quelques millivolts.
    \item \'Etendue de mesure : $T_{ARM} \in [-20\,;\,200]\ \si{\celsius}$, soit une \'etendue de $\SI{220}{\celsius}$.
    \item Tensions extr\^emes : $e_{AB\,mini} = \SI{-1}{\milli\volt}$ (\`a $\SI{-20}{\celsius}$) et $e_{AB\,maxi} = \SI{10,8}{\milli\volt}$ (\`a $\SI{200}{\celsius}$), soit une dynamique utile de $\SI{11,8}{\milli\volt}$ seulement.
  \end{itemize}
  \item Application de la relation fournie, en convertissant les millivolts en volts :
  \[ V_4(e_{AB}{=}\SI{-1}{\milli\volt}) = 337{,}7\times(-0{,}001) + 1{,}34 = -0{,}3377 + 1{,}34 \approx \SI{1}{\volt} \]
  \[ V_4(e_{AB}{=}\SI{10,8}{\milli\volt}) = 337{,}7\times 0{,}0108 + 1{,}34 = 3{,}647 + 1{,}34 \approx \SI{5}{\volt} \]
  \item Le profil \'etant lin\'eaire (hypoth\`ese explicite du DT10), la sensibilit\'e est constante :
  \[ \dfrac{\Delta V_4}{\Delta T_{ARM}} = \dfrac{5-1}{200-(-20)} = \dfrac{4}{220} \approx \SI{18,2}{\milli\volt\per\celsius} \]
  \item Le terme $1{,}34$ est la tension de \textbf{d\'ecalage} $V_d$ apport\'ee par le deuxi\`eme \'etage. Sans lui, on aurait $V_4(\SI{-20}{\celsius}) = 337{,}7\times(-0{,}001) = \SI{-0,338}{\volt}$ : une tension \textbf{n\'egative}, que l'\'etage de conversion tension-courant qui suit ne peut pas transformer en un courant positif, et qui sortirait de la plage $1$--$\SI{5}{\volt}$ attendue en entr\'ee de cet \'etage. Le d\'ecalage translate donc toute la caract\'eristique vers les tensions positives.
\end{enumerate}

\textbf{Conclusion.} Une dynamique de capteur de $\SI{11,8}{\milli\volt}$ seulement est transform\'ee, par amplification puis d\'ecalage, en la fen\^etre $1$--$\SI{5}{\volt}$ exactement : c'est la condition pour que l'\'etage suivant d\'elivre la plage normalis\'ee $4$--$\SI{20}{\milli\ampere}$ annonc\'ee par le DT10, courant de repos de $\SI{4}{\milli\ampere}$ compris. Sans lecture pr\'ealable du document technique (nature du capteur, bornes de la plage), aucun des deux \'etages ne peut \^etre dimensionn\'e.
}
\end{tcolorbox}

\begin{tcolorbox}[exoBox,
  title={Exercice 3 -- \difficulte{2} \hfill\normalfont\footnotesize 2024, Extracteur d'air, partie E, question E4}
]
\CAPETSHOW{
La tension $V_3$ obtenue en sortie de l'\'etage d'amplification du transmetteur fait office d'entr\'ee $V_e$ pour un \'etage de d\'ecalage en tension. Ce montage, construit autour d'un AOP suppos\'e id\'eal, produit une tension $V_4$ qui d\'ecale $V_e$ d'une valeur de r\'ef\'erence not\'ee $V_3'$ (notation locale, \`a ne pas confondre avec le $V_3$ de l'\'etage d'amplification). L'entr\'ee non-inverseuse re\c coit les tensions $V_3'$ et $V_e$ chacune \`a travers une r\'esistance $R$, l'entr\'ee inverseuse est reboucl\'ee sur la sortie $V_4$ \`a travers une r\'esistance $R$ et reli\'ee \`a la masse \`a travers une r\'esistance $R$ \'egalement.

\begin{center}
\begin{circuitikz}[european, scale=0.85]
\draw (4,0) node[op amp, yscale=-1] (U) {};
\draw (U.+) -- ++(-1,0) coordinate(P) node[circ]{};
\node[above right=1pt and 2pt] at (P) {$V_+$};
\draw (P) ++(0,0.6) coordinate(P1) -- ++(-2,0) to[R, l=$R$] ++(-2,0) node[left]{$V_3'$};
\draw (P) ++(0,-0.6) coordinate(P2) -- ++(-2,0) to[R, l=$R$] ++(-2,0) node[left]{$V_e$};
\draw (P1) -- (P2);
\draw (U.out) -- ++(1,0) coordinate(S) node[right]{$V_4$};
\draw (U.-) -- ++(-2.6,0) coordinate(N);
\node[below left=1pt and 1pt] at (N) {$V_-$};
\draw (S) to[R, l=$R$] ++(0,-2) coordinate(M) node[circ]{};
\draw (N) -- (N |- M) -- (M);
\draw (M) to[R, l=$R$] ++(0,-1.6) node[ground]{};
\end{circuitikz}
\end{center}

\begin{enumerate}[leftmargin=1.5em]
  \item En appliquant le th\'eor\`eme de Millman au n\oe ud $V_+$ (deux r\'esistances $R$, l'AOP ne tirant aucun courant), exprimer $V_+$ en fonction de $V_3'$ et $V_e$.
  \item En appliquant le th\'eor\`eme de Millman (ou le diviseur de tension) au n\oe ud $V_-$, exprimer $V_-$ en fonction de $V_4$.
  \item Montrer qu'en r\'egime lin\'eaire, $V_4 = V_3' + V_e$.
\end{enumerate}

}{
\textbf{Corrig\'e.}

\begin{enumerate}[leftmargin=1.5em]
  \item Le n\oe ud $V_+$ re\c coit deux branches de m\^eme r\'esistance $R$ (depuis $V_3'$ et depuis $V_e$), et aucun courant ne p\'en\`etre dans l'AOP par $V_+$ :
  \[ V_+ = \dfrac{\dfrac{V_3'}{R} + \dfrac{V_e}{R}}{\dfrac{1}{R}+\dfrac{1}{R}} = \dfrac{V_3' + V_e}{2} \]
  \item De m\^eme, $V_-$ est le point milieu d'un diviseur \`a deux r\'esistances \'egales $R$ entre $V_4$ (sortie) et la masse :
  \[ V_- = \dfrac{\dfrac{V_4}{R} + \dfrac{0}{R}}{\dfrac{1}{R}+\dfrac{1}{R}} = \dfrac{V_4}{2} \]
  \item En r\'egime lin\'eaire (contre-r\'eaction n\'egative sur $V_-$), $V_+ = V_-$ :
  \[ \dfrac{V_3'+V_e}{2} = \dfrac{V_4}{2} \quad\Longrightarrow\quad V_4 = V_3' + V_e \]
\end{enumerate}

\textbf{Conclusion.} L'\'etage r\'ealise une addition exacte, \`a gain unitaire, du signal amplifi\'e $V_e$ et de la r\'ef\'erence $V_3'$ : c'est ce d\'ecalage qui translate la sortie de l'amplificateur vers la fen\^etre $1$--$\SI{5}{\volt}$ exig\'ee \`a l'entr\'ee du convertisseur tension-courant. Sans lui, la temp\'erature minimale du capteur correspondrait \`a $\SI{0}{\milli\ampere}$, indiscernable d'une rupture de la ligne de transmission, alors que la boucle normalis\'ee impose un courant de repos de $\SI{4}{\milli\ampere}$.
}
\end{tcolorbox}

\begin{tcolorbox}[exoBox,
  title={Exercice 4 -- \difficulte{2} \hfill\normalfont\footnotesize exercice in\'edit}
]
\CAPETSHOW{
Pour chacun des besoins fonctionnels list\'es dans le tableau ci-dessous, identifier le montage \`a AOP le plus adapt\'e parmi : suiveur, amplificateur inverseur, amplificateur non-inverseur, soustracteur (diff\'erentiel), comparateur \`a hyst\'er\'esis.

\begin{center}
\begin{tabular}{p{8.5cm}p{4.5cm}}
\toprule
\textbf{Besoin fonctionnel} & \textbf{Montage \`a choisir} \\
\midrule
Isoler un capteur haute imp\'edance sans att\'enuer ni inverser le signal & \dotfill \\
Amplifier une tension diff\'erentielle issue d'un pont de mesure (deux bornes flottantes) & \dotfill \\
Basculer une sortie logique en \'evitant les oscillations autour d'un seuil bruit\'e & \dotfill \\
Amplifier un signal en inversant sa phase, avec une imp\'edance d'entr\'ee fix\'ee et connue & \dotfill \\
Amplifier un signal en conservant sa phase, avec une imp\'edance d'entr\'ee tr\`es \'elev\'ee & \dotfill \\
\bottomrule
\end{tabular}
\end{center}

Pour chaque r\'eponse, justifier bri\`evement le choix en une phrase.

}{
\textbf{Corrig\'e.}

\begin{center}
\begin{tabular}{p{8.5cm}p{4.5cm}}
\toprule
\textbf{Besoin fonctionnel} & \textbf{Montage} \\
\midrule
Isoler un capteur haute imp\'edance sans att\'enuer ni inverser le signal & Suiveur \\
Amplifier une tension diff\'erentielle issue d'un pont de mesure (deux bornes flottantes) & Soustracteur \\
Basculer une sortie logique en \'evitant les oscillations autour d'un seuil bruit\'e & Comparateur \`a hyst\'er\'esis \\
Amplifier un signal en inversant sa phase, avec une imp\'edance d'entr\'ee fix\'ee et connue & Amplificateur inverseur \\
Amplifier un signal en conservant sa phase, avec une imp\'edance d'entr\'ee tr\`es \'elev\'ee & Amplificateur non-inverseur \\
\bottomrule
\end{tabular}
\end{center}

\begin{itemize}
  \item Le \textbf{suiveur} ($V_s=V_e$, gain unitaire) a pour seul r\^ole l'adaptation d'imp\'edance : imp\'edance d'entr\'ee infinie, imp\'edance de sortie nulle, aucune att\'enuation ni inversion.
  \item Le \textbf{soustracteur} calcule directement $V_s = (R_2/R_1)\,(V_2-V_1)$ : c'est le seul montage de base qui traite nativement deux entr\'ees flottantes (aucune des deux n'est r\'ef\'erenc\'ee \`a la masse), typique d'un pont de Wheatstone.
  \item Le \textbf{comparateur \`a hyst\'er\'esis} introduit deux seuils distincts (contre-r\'eaction positive) : contrairement \`a un comparateur simple \`a seuil unique, un bruit de faible amplitude autour du seuil ne provoque plus de basculements r\'ep\'et\'es intempestifs.
  \item L'\textbf{amplificateur inverseur} a une imp\'edance d'entr\'ee \'egale \`a $R_1$ (fix\'ee par le concepteur), ce qui est parfois recherch\'e (par exemple pour sommer plusieurs signaux sur un m\^eme n\oe ud avec des poids connus), au prix d'une inversion de phase.
  \item L'\textbf{amplificateur non-inverseur} pr\'esente une imp\'edance d'entr\'ee quasi infinie (entr\'ee directement sur $V_+$) tout en amplifiant (gain $\geq 1$) sans inverser la phase, contrairement \`a l'inverseur.
\end{itemize}

\textbf{Conclusion.} Le choix d'un montage ne se d\'eduit pas d'une formule mais du couple \guillemotleft\ imp\'edance vue par la source / traitement attendu du signal \guillemotright\ : c'est ce raisonnement, et non le seul calcul de gain, qui est \'evalu\'e \`a l'\'ecrit lorsqu'une question demande de justifier une structure de conditionnement.
}
\end{tcolorbox}

\begin{tcolorbox}[exoBox,
  title={Exercice 5 -- \difficulte{2} \hfill\normalfont\footnotesize 2026, \'epreuve BAC3, questions B7 et B8 (document technique DTB3)}
]
\CAPETSHOW{
Poursuite de l'\'etude de la cha\^ine d'acquisition de temp\'erature de la pile \`a combustible (exercice 1). En sortie du conditionneur du capteur, la tension disponible est $v_T = U_0\,(1+a\,T)$, avec $U_0 = \SI{1}{\volt}$, qui \'evolue donc autour de $\SI{1}{\volt}$ avec une faible dynamique. Le convertisseur analogique-num\'erique plac\'e en fin de cha\^ine impose en revanche, \`a l'entr\'ee, une tension $v_A$ \emph{directement proportionnelle} \`a la temp\'erature, $v_A = c\cdot T$ avec $c>0$, \'evoluant sur une plage de $\SI{3}{\volt}$. On doit donc ins\'erer un ensemble \{amplificateur + conditionneur\} construit \`a partir des montages du document technique ci-dessous.

\begin{center}
\footnotesize
\textbf{\itshape Document technique DTB3 -- Montages \`a ALI disponibles (ALI id\'eaux, $\varepsilon = 0$)}
\smallskip
\renewcommand{\arraystretch}{1.5}
\begin{tabular}{@{}llp{5.4cm}@{}}
\toprule
\textbf{Montage} & \textbf{Fonction de transfert} & \textbf{C\^ablage} \\
\midrule
Montage 1 & $\dfrac{v_S}{v_E} = \dfrac{R_1+R_0}{R_1}$ & $v_E$ sur l'entr\'ee $+$ ; pont $R_0$ (contre-r\'eaction) -- $R_1$ (vers la masse) sur l'entr\'ee $-$ \\
Montage 2 & $\dfrac{v_S}{v_E} = -\dfrac{R_0}{R_1}$ & $v_E$ via $R_1$ sur l'entr\'ee $-$ ; $R_0$ en contre-r\'eaction ; entr\'ee $+$ \`a la masse \\
Montage 3 & $v_S = \dfrac{R_0}{R_1}\,(v_2 - v_1)$ & deux entr\'ees $v_1$ et $v_2$, quatre r\'esistances appari\'ees $R_1$ et $R_0$ \\
Montage 4 & (quatre r\'esistances $R_1$, $R_2$, $R_3$, $R_0$, deux entr\'ees $v_1$, $v_2$) & structure diff\'erentielle \`a r\'esistances non appari\'ees \\
\bottomrule
\end{tabular}
\end{center}

\begin{enumerate}[leftmargin=1.5em]
  \item \textbf{D\'emontrer} la fonction de transfert du montage 1 \`a partir des hypoth\`eses de l'ALI id\'eal. En d\'eduire le nom et le r\^ole de ce type de montage.
  \item Le montage 1 permet-il, \`a lui seul, d'obtenir $v_A = c\cdot T$ \`a partir de $v_T = U_0(1+aT)$~? Justifier par le calcul de $v_A$ qu'il produirait.
  \item Choisir, en justifiant, le montage du document technique DTB3 qui permet de r\'ealiser l'ensemble \{amplificateur + conditionneur\}. Pr\'eciser ce qu'il faut appliquer sur chacune de ses entr\'ees.
\end{enumerate}

}{
\textbf{Corrig\'e.}

\begin{enumerate}[leftmargin=1.5em]
  \item \textbf{Montage 1 :} $\dfrac{v_S}{v_E} = \dfrac{R_1+R_0}{R_1}$.

  \textbf{D\'emonstration.} La sortie est reboucl\'ee sur l'entr\'ee $-$ : contre-r\'eaction n\'egative, donc r\'egime lin\'eaire et $V_+ = V_-$. Or $V_+ = v_E$ (branchement direct). Comme $I_- = 0$, $R_0$ et $R_1$ forment un simple diviseur de tension non charg\'e entre $v_S$ et la masse :
  \[ V_- = v_S\,\dfrac{R_1}{R_1+R_0} \]
  L'\'egalit\'e $V_+ = V_-$ donne $v_E = v_S\dfrac{R_1}{R_1+R_0}$, soit
  \[ \dfrac{v_S}{v_E} = \dfrac{R_1+R_0}{R_1} = 1 + \dfrac{R_0}{R_1} \]
  C'est l'\textbf{amplificateur non-inverseur} : il amplifie d'un gain toujours $\geq 1$ sans inverser la phase, en pr\'esentant une imp\'edance d'entr\'ee quasi infinie (entr\'ee directement sur $V_+$) -- il ne charge donc pas l'\'etage amont.
  \item \textbf{Non.} En posant $G = 1+R_0/R_1 > 0$, ce montage donnerait
  \[ v_A = G\cdot v_T = G\,U_0(1+aT) = \underbrace{G\,U_0}_{\text{terme constant}} + \underbrace{G\,U_0\,a}_{\text{pente}}\;T \]
  La sortie contient un \textbf{terme constant non nul} $G\,U_0$ (l'\og\ offset\ \fg{} de la sonde \`a $\SI{0}{\celsius}$), alors que le cahier des charges exige $v_A = c\cdot T$, c'est-\`a-dire une caract\'eristique passant par l'origine. De plus, $G$ ne peut pas \^etre choisi librement : le rendre assez grand pour couvrir $\SI{3}{\volt}$ de dynamique utile amplifierait d'autant le terme constant et saturerait l'ALI. Un montage \`a une seule entr\'ee ne peut donc pas r\'ealiser simultan\'ement l'amplification \emph{et} la soustraction de l'offset.
  \item Il faut le \textbf{montage 3}, le soustracteur (amplificateur diff\'erentiel \`a r\'esistances appari\'ees), de fonction de transfert $v_S = \dfrac{R_0}{R_1}(v_2 - v_1)$. Il r\'ealise en un seul \'etage les deux fonctions attendues :
  \begin{itemize}
    \item on applique sur $v_2$ la tension issue du conditionneur, $v_T$ ;
    \item on applique sur $v_1$ une \textbf{source de tension de r\'ef\'erence constante} \'egale \`a $U_0$ (la valeur de $v_T$ \`a $\SI{0}{\celsius}$), ce qui \'elimine exactement le terme constant : $v_2-v_1 = U_0(1+aT) - U_0 = U_0\,a\,T$ ;
    \item le rapport $R_0/R_1$ assure ensuite l'\textbf{amplification} : $v_A = \dfrac{R_0}{R_1}U_0\,a\,T$, qui est bien de la forme $c\cdot T$ avec $c>0$.
  \end{itemize}
  Le montage 2 est \'ecart\'e (une seule entr\'ee, et surtout gain n\'egatif incompatible avec $c>0$), le montage 1 par la question pr\'ec\'edente, et le montage 4 -- diff\'erentiel \`a r\'esistances non appari\'ees -- n'apporte ici aucun degr\'e de libert\'e utile tout en d\'egradant la r\'ejection de mode commun.
\end{enumerate}

\textbf{Conclusion.} Le choix ne se fait pas sur la seule valeur de gain, mais sur la \emph{forme} de la relation entr\'ee-sortie exig\'ee : seule une structure diff\'erentielle permet d'annuler l'offset $U_0$ de la sonde tout en amplifiant sa variation, condition n\'ecessaire pour que la tension $v_A = c\,T$ balaie toute la plage de $\SI{3}{\volt}$ du convertisseur et que la r\'esolution de la mesure -- donc l'exigence d'erreur inf\'erieure \`a $\SI{1}{\celsius}$ -- soit tenue.
}
\end{tcolorbox}

\begin{tcolorbox}[exoBox,
  title={Exercice 6 -- \difficulte{3} \hfill\normalfont\footnotesize 2023, \'epreuve \'ecrite disciplinaire, questions F9 et F10 (dimensionnement)}
]
\CAPETSHOW{
Le robot enjambeur Bakus \'evolue dans des vignes de coteau. Deux centrales inertielles mesurent son assiette : leurs sorties analogiques, une fois trait\'ees, \'evoluent \textbf{proportionnellement} entre $\SI{0}{\volt}$ et $\SI{5}{\volt}$ pour des inclinaisons comprises entre $0\,\%$ et $+50\,\%$ de pente. Cette tension constitue l'entr\'ee $V_{e2}$ du montage ci-dessous.

\smallskip
Le cahier des charges de s\'ecurit\'e impose deux basculements \emph{distincts} : le robot passe du mode autonome au mode manuel d\`es qu'une pente sup\'erieure \`a $\mathbf{25\,\%}$ est d\'etect\'ee, et le retour au mode autonome n'est possible qu'apr\`es intervention humaine \emph{et} retour sur un terrain de pente inf\'erieure \`a $\mathbf{23\,\%}$. On r\'ealise cette fonction par un comparateur inverseur \`a hyst\'er\'esis \`a ALI id\'eal, dont les tensions de saturation en sortie valent $\pm\SI{15}{\volt}$. On rappelle que pour ce montage
\[ V_+ = \frac{R_2\,V_{ref2} + R_1\,V_{s2}}{R_1+R_2} \]

\begin{center}
\begin{circuitikz}[european, scale=0.85]
\draw (4,0) node[op amp, yscale=-1] (U) {};
\draw (U.-) -- ++(-1.6,0) node[left]{$V_{e2}$};
\draw (U.+) -- ++(-1,0) coordinate(P) node[circ]{};
\node[above right=1pt and 2pt] at (P) {$V_+$};
\draw (P) to[R, l=$R_1$] ++(-2.6,0) coordinate(REF);
\draw (REF) to[V, l_=$V_{ref2}$] ++(0,-2.6) node[ground]{};
\draw (U.out) -- ++(1,0) coordinate(S) node[circ]{} -- ++(0.9,0) node[right]{$V_{s2}$};
\draw (S) to[R, l_=$R_2$] ++(0,2.2) coordinate(F) -- (F -| P) -- (P);
\end{circuitikz}
\end{center}

\begin{enumerate}[leftmargin=1.5em]
  \item \'Etablir la relation lin\'eaire liant la tension $V_{e2}$ \`a la pente $p$ (en \%), puis en d\'eduire les deux tensions de seuil $V_H$ et $V_B$ qu'impose le cahier des charges de s\'ecurit\'e.
  \item En partant de l'expression rappel\'ee de $V_+$, exprimer la \textbf{largeur d'hyst\'er\'esis} $V_H - V_B$ en fonction de $R_1$, $R_2$ et $V_{sat}$. En prenant $R_1 = \SI{1}{\kilo\ohm}$, en d\'eduire la valeur de $R_2$ n\'ecessaire.
  \item Exprimer la \textbf{valeur moyenne} $(V_H+V_B)/2$ des deux seuils en fonction de $R_1$, $R_2$ et $V_{ref2}$, et en d\'eduire la valeur que doit prendre la source $V_{ref2}$.
  \item Le concepteur retient finalement les valeurs normalis\'ees $R_2 = \SI{150}{\kilo\ohm}$ et $V_{ref2} = \SI{2,42}{\volt}$. V\'erifier que le montage reste conforme au cahier des charges, puis tracer l'allure de la caract\'eristique $V_{s2} = f(V_{e2})$ en indiquant le sens de parcours.
\end{enumerate}

}{
\textbf{Corrig\'e.}

\begin{enumerate}[leftmargin=1.5em]
  \item La proportionnalit\'e $0$--$\SI{5}{\volt}$ pour $0$--$50\,\%$ donne une pente de conversion de $\SI{0,1}{\volt}$ par point de pente :
  \[ V_{e2} = \dfrac{5}{50}\,p = 0{,}1\,p \quad (p \text{ en \%}) \]
  Le basculement vers le mode manuel se produit \`a $p = 25\,\%$ et le retour au mode autonome \`a $p = 23\,\%$, d'o\`u les deux seuils vis\'es :
  \[ V_H = 0{,}1\times 25 = \SI{2,5}{\volt} \qquad V_B = 0{,}1\times 23 = \SI{2,3}{\volt} \]
  \item Les deux seuils correspondent aux deux \'etats de saturation de la sortie, $V_{s2} = \pm V_{sat}$ :
  \[ V_H = \dfrac{R_2 V_{ref2} + R_1 V_{sat}}{R_1+R_2} \qquad V_B = \dfrac{R_2 V_{ref2} - R_1 V_{sat}}{R_1+R_2} \]
  Par diff\'erence, le terme en $V_{ref2}$ dispara\^it :
  \[ V_H - V_B = \dfrac{2\,R_1\,V_{sat}}{R_1+R_2} \]
  La largeur d'hyst\'er\'esis ne d\'epend donc que du rapport des r\'esistances et de $V_{sat}$. En imposant $V_H-V_B = 2{,}5-2{,}3 = \SI{0,2}{\volt}$ avec $V_{sat}=\SI{15}{\volt}$ :
  \[ \dfrac{R_1}{R_1+R_2} = \dfrac{0{,}2}{2\times 15} = \dfrac{1}{150} \quad\Longrightarrow\quad R_2 = 149\,R_1 = \SI{149}{\kilo\ohm} \]
  \item Par demi-somme, cette fois c'est le terme en $V_{sat}$ qui dispara\^it :
  \[ \dfrac{V_H+V_B}{2} = \dfrac{R_2\,V_{ref2}}{R_1+R_2} \]
  Le centre de l'hyst\'er\'esis vis\'e vaut $(2{,}5+2{,}3)/2 = \SI{2,4}{\volt}$, d'o\`u :
  \[ V_{ref2} = \SI{2,4}{\volt}\times\dfrac{R_1+R_2}{R_2} = 2{,}4\times\dfrac{150\,000}{149\,000} \approx \SI{2,416}{\volt} \]
  \item Avec les valeurs normalis\'ees $R_1 = \SI{1}{\kilo\ohm}$, $R_2 = \SI{150}{\kilo\ohm}$ et $V_{ref2} = \SI{2,42}{\volt}$ :
  \[ V_H = \dfrac{150\,000\times 2{,}42 + 1\,000\times 15}{151\,000} = \dfrac{378\,000}{151\,000} \approx \SI{2,503}{\volt} \]
  \[ V_B = \dfrac{150\,000\times 2{,}42 - 1\,000\times 15}{151\,000} = \dfrac{348\,000}{151\,000} \approx \SI{2,305}{\volt} \]
  Ramen\'es en pente par $p = V/0{,}1$ : $p_H \approx 25{,}0\,\%$ et $p_B \approx 23{,}0\,\%$. Les deux exigences sont donc tenues, l'\'ecart aux valeurs th\'eoriques ($+0{,}03$ point de pente) \'etant tr\`es inf\'erieur \`a la dispersion des composants.

  \textbf{Caract\'eristique} (comparateur \emph{inverseur} : $V_{e2}$ attaque l'entr\'ee \guillemotleft\ -- \guillemotright, donc $V_{s2}$ \textbf{d\'ecro\^it} quand $V_{e2}$ cro\^it -- caract\'eristique descendante) :
  \begin{center}
  \begin{tikzpicture}[scale=0.95, font=\footnotesize]
    \draw[->] (0,0) -- (5.6,0) node[right] {$V_{e2}$ (V)};
    \draw[->] (0,-1.9) -- (0,1.9) node[above] {$V_{s2}$ (V)};
    \draw[thick, colExo] (0.3,1.3) -- (3.4,1.3);
    \draw[thick, colExo, ->] (3.4,1.3) -- (3.4,-1.3);
    \draw[thick, colExo] (3.4,-1.3) -- (5.2,-1.3);
    \draw[thick, colExo!55] (1.6,-1.3) -- (3.4,-1.3);
    \draw[thick, colExo!55, ->] (1.6,-1.3) -- (1.6,1.3);
    \draw[thick, colExo!55] (0.3,1.3) -- (1.6,1.3);
    \draw[dashed] (3.4,0) -- (3.4,1.3);
    \draw[dashed] (1.6,0) -- (1.6,-1.3);
    \node[below right] at (3.4,0) {\scriptsize $V_H = 2{,}50$};
    \node[above left] at (1.6,0) {\scriptsize $V_B = 2{,}31$};
    \node[left] at (0,1.3) {\scriptsize $+15$};
    \node[left] at (0,-1.3) {\scriptsize $-15$};
  \end{tikzpicture}
  \end{center}
  Trajet fonc\'e : $V_{e2}$ croissant (pente augmentant), basculement vers $-V_{sat}$ en $V_H$ -- passage en mode manuel. Trajet clair : $V_{e2}$ d\'ecroissant, retour \`a $+V_{sat}$ seulement en $V_B$.
\end{enumerate}

\textbf{Conclusion.} La d\'emarche de dimensionnement se lit directement sur les deux combinaisons des seuils : la \textbf{diff\'erence} $V_H-V_B$ ne fixe que le rapport $R_1/R_2$, la \textbf{demi-somme} ne fixe que $V_{ref2}$ -- les deux exigences du cahier des charges (basculement \`a $25\,\%$, retour \`a $23\,\%$) sont donc r\'eglables ind\'ependamment. C'est cet \'ecart de 2 points de pente, et non un seuil unique, qui interdit au robot de basculer sans fin entre mode autonome et mode manuel lorsqu'il roule au voisinage de la pente critique.
}
\end{tcolorbox}

\CAPETENDFILE
